\documentclass{amsart}
\usepackage[UTF8,fontset=windows]{ctex}
\usepackage[a4paper,margin=24mm]{geometry}
% For dealing with references we use the comment environment
\usepackage{verbatim}
\newenvironment{reference}{\comment}{\endcomment}
%\newenvironment{reference}{}{}
\newenvironment{slogan}{\comment}{\endcomment}
\newenvironment{history}{\comment}{\endcomment}

% For commutative diagrams we use Xy-pic
\usepackage[all]{xy}

% We use 2cell for 2-commutative diagrams.
\xyoption{2cell}
\UseAllTwocells

% We use multicol for the list of chapters between chapters
\usepackage{multicol}

% This is generally recommended for better output
\usepackage{lmodern}
\usepackage[T1]{fontenc}

% For cross-file-references

% Package for hypertext links:
\usepackage{hyperref}

% For any local file, say "hello.tex" you want to link to please

% Theorem environments.
%
\theoremstyle{plain}
\newtheorem{theorem}[subsection]{Theorem}
\newtheorem{proposition}[subsection]{Proposition}
\newtheorem{lemma}[subsection]{Lemma}

\theoremstyle{definition}
\newtheorem{definition}[subsection]{Definition}
\newtheorem{example}[subsection]{Example}
\newtheorem{exercise}[subsection]{Exercise}
\newtheorem{situation}[subsection]{Situation}

\theoremstyle{remark}
\newtheorem{remark}[subsection]{Remark}
\newtheorem{remarks}[subsection]{Remarks}

\numberwithin{equation}{subsection}

% Macros
%
\def\lim{\mathop{\mathrm{lim}}\nolimits}
\def\colim{\mathop{\mathrm{colim}}\nolimits}
\def\Spec{\mathop{\mathrm{Spec}}}
\def\Hom{\mathop{\mathrm{Hom}}\nolimits}
\def\Ext{\mathop{\mathrm{Ext}}\nolimits}
\def\SheafHom{\mathop{\mathcal{H}\!\mathit{om}}\nolimits}
\def\SheafExt{\mathop{\mathcal{E}\!\mathit{xt}}\nolimits}
\def\Sch{\mathit{Sch}}
\def\Mor{\mathop{\mathrm{Mor}}\nolimits}
\def\Ob{\mathop{\mathrm{Ob}}\nolimits}
\def\Sh{\mathop{\mathit{Sh}}\nolimits}
\def\NL{\mathop{N\!L}\nolimits}
\def\CH{\mathop{\mathrm{CH}}\nolimits}
\def\proetale{{pro\text{-}\acute{e}tale}}
\def\etale{{\acute{e}tale}}
\def\QCoh{\mathit{QCoh}}
\def\Ker{\mathop{\mathrm{Ker}}}
\def\Im{\mathop{\mathrm{Im}}}
\def\Coker{\mathop{\mathrm{Coker}}}
\def\Coim{\mathop{\mathrm{Coim}}}

% Boxtimes
%
\DeclareMathSymbol{\boxtimes}{\mathbin}{AMSa}{"02}

%
% Macros for moduli stacks/spaces
%
\def\QCohstack{\mathcal{QC}\!\mathit{oh}}
\def\Cohstack{\mathcal{C}\!\mathit{oh}}
\def\Spacesstack{\mathcal{S}\!\mathit{paces}}
\def\Quotfunctor{\mathrm{Quot}}
\def\Hilbfunctor{\mathrm{Hilb}}
\def\Curvesstack{\mathcal{C}\!\mathit{urves}}
\def\Polarizedstack{\mathcal{P}\!\mathit{olarized}}
\def\Complexesstack{\mathcal{C}\!\mathit{omplexes}}
% \Pic is the operator that assigns to X its picard group, usage \Pic(X)
% \Picardstack_{X/B} denotes the Picard stack of X over B
% \Picardfunctor_{X/B} denotes the Picard functor of X over B
\def\Pic{\mathop{\mathrm{Pic}}\nolimits}
\def\Picardstack{\mathcal{P}\!\mathit{ic}}
\def\Picardfunctor{\mathrm{Pic}}
\def\Deformationcategory{\mathcal{D}\!\mathit{ef}}

\hypersetup{hidelinks}
\widowpenalty=10000\clubpenalty=10000\displaywidowpenalty=10000
\begin{document}
\title[Stacks Project, Tag 0DUT]{有限惯性代数栈的模空间装配}
\maketitle
\noindent\textbf{Counting scope.} The primary problem is Tag 0DUT, marked below. All preceding numbered statements are supporting author context and are not counted as additional problems. This is a local pilot with a temporary ID, not production admission.


\noindent Copyright (C) 2005--2025 Johan de Jong. Stacks Project authors. Source: \href{https://stacks.math.columbia.edu/tag/0DUT}{Tag 0DUT}.

\noindent Difficulty H3: 在给定étale局部有限惯性覆盖及well-nigh-affine局部模空间后，仍须构造整体模空间与模射。证明自行组成M\_ij的étale等价关系、验证几何点上的关系并粘合f，再核实基本图表为笛卡尔及任意平坦基变换后的范畴泛性质。局部存在性是被引输入，整体存在性与泛性质装配仍有独立实质内容，因此列为H3候选。.

\noindent History: excerpt from revision local cached snapshot; exact chapter hashes; revision unknown, stacks-more-morphisms.tex, lines 3636--3885. Only reference presentation was adapted. Licensed under GNU FDL 1.2 or later; no invariant sections or cover texts. See accompanying COPYING.

\begin{lemma}
\label{lemma-etale-local-finite-inertia}
Let $\mathcal{X}$ be an algebraic stack. Assume
$\mathcal{I}_\mathcal{X} \to \mathcal{X}$ is finite.
Then there exist a set $I$ and for $i \in I$ a morphism of algebraic stacks
$$
g_i : \mathcal{X}_i \longrightarrow \mathcal{X}
$$
with the following properties
\begin{enumerate}
\item $|\mathcal{X}| = \bigcup |g_i|(|\mathcal{X}_i|)$,
\item $\mathcal{X}_i$ is well-nigh affine,
\item $\mathcal{I}_{\mathcal{X}_i} \to
\mathcal{X}_i \times_\mathcal{X} \mathcal{I}_\mathcal{X}$
is an isomorphism, and
\item $g_i : \mathcal{X}_i \to \mathcal{X}$ is representable
by algebraic spaces, separated, and \'etale,
\end{enumerate}
\end{lemma}

\begin{proof}
For any $x \in |\mathcal{X}|$ we can choose
$g : \mathcal{U} \to \mathcal{X}$, $\mathcal{U} = [U/R]$, and $u$ as in
Morphisms of Stacks,
Lemma \href{https://stacks.math.columbia.edu/tag/0DU4}{lemma etale local quasi DM at x (Tag 0DU4)}.
Then by
Morphisms of Stacks, Lemma
\href{https://stacks.math.columbia.edu/tag/0DUA}{lemma stabilizer preserving unramified (Tag 0DUA)}
we see that there exists an open substack
$\mathcal{U}' \subset \mathcal{U}$ containing $u$
such that $\mathcal{I}_{\mathcal{U}'} \to
\mathcal{U}' \times_\mathcal{X} \mathcal{I}_\mathcal{X}$
is an isomorphism.
Let $U' \subset U$ be the $R$-invariant open corresponding to
the open substack $\mathcal{U}'$.
Let $u' \in U'$ be a point of $U'$ mapping to $u$.
Observe that $t(s^{-1}(\{u'\}))$ is finite as $s : R \to U$ is finite.
By Properties, Lemma \href{https://stacks.math.columbia.edu/tag/01ZY}{lemma ample finite set in affine (Tag 01ZY)}
and Groupoids, Lemma \href{https://stacks.math.columbia.edu/tag/03JE}{lemma find invariant affine (Tag 03JE)}
we can find an $R$-invariant affine open $U'' \subset U'$
containing $u'$. Let $R''$ be the restriction of $R$ to $U''$.
Then $\mathcal{U}'' = [U''/R'']$ is an open substack of
$\mathcal{U}'$ containing $u$, is well-nigh affine,
$\mathcal{I}_{\mathcal{U}''} \to
\mathcal{U}'' \times_\mathcal{X} \mathcal{I}_\mathcal{X}$
is an isomorphism, and $\mathcal{U}'' \to \mathcal{X}$
and is representable by algebraic spaces and \'etale.
Finally, $\mathcal{U}'' \to \mathcal{X}$ is separated as
$\mathcal{U}''$ is separated (Lemma \href{https://stacks.math.columbia.edu/tag/0DUM}{lemma well nigh affine (Tag 0DUM)})
the diagonal of $\mathcal{X}$ is separated
(Morphisms of Stacks, Lemma \href{https://stacks.math.columbia.edu/tag/0CL0}{lemma diagonal diagonal (Tag 0CL0)})
and separatedness follows from Morphisms of Stacks, Lemma
\href{https://stacks.math.columbia.edu/tag/050M}{lemma compose after separated (Tag 050M)}.
Since the point $x \in |\mathcal{X}|$ is arbitrary the proof is complete.
\end{proof}


\begin{lemma}
\label{lemma-well-nigh-affine-moduli-space}
Let the algebraic stack $\mathcal{X}$ be well-nigh affine. There exists
a uniform categorical moduli space
$$
f : \mathcal{X} \longrightarrow M
$$
in the category of affine schemes. Moreover
$f$ is separated, quasi-compact, and a universal homeomorphism.
\end{lemma}

\begin{proof}
Write $\mathcal{X} = [U/R]$ with $(U, R, s, t, c)$ as in
Lemma \href{https://stacks.math.columbia.edu/tag/0DUM}{lemma well nigh affine (Tag 0DUM)}. Let $C$ be the ring of
$R$-invariant functions on $U$, see
Groupoids, Section \href{https://stacks.math.columbia.edu/tag/03BE}{section finite flat (Tag 03BE)}.
We set $M = \Spec(C)$. The $R$-invariant morphism
$U \to M$ corresponds to a morphism $f : \mathcal{X} \to M$ by
Lemma \href{https://stacks.math.columbia.edu/tag/0DUH}{lemma quotient compare (Tag 0DUH)}.
The characterization of morphisms into affine schemes given in
Schemes, Lemma \href{https://stacks.math.columbia.edu/tag/01I1}{lemma morphism into affine (Tag 01I1)}
immediately guarantees that $\phi : U \to M$ is a categorical
quotient in the category of affine schemes. Hence $f$ is a
categorical moduli space in the category of affine schemes
(Lemma \href{https://stacks.math.columbia.edu/tag/0DUI}{lemma categorical quotient compare (Tag 0DUI)}).

\medskip\noindent
Since $\mathcal{X}$ is separated by Lemma \href{https://stacks.math.columbia.edu/tag/0DUM}{lemma well nigh affine (Tag 0DUM)}
we find that $f$ is separated by Morphisms of Stacks, Lemma
\href{https://stacks.math.columbia.edu/tag/050M}{lemma compose after separated (Tag 050M)}.

\medskip\noindent
Since $U \to \mathcal{X}$ is surjective and since $U \to M$ is quasi-compact,
we see that $f$ is quasi-compact by Morphisms of Stacks, Lemma
\href{https://stacks.math.columbia.edu/tag/050X}{lemma surjection from quasi compact (Tag 050X)}.

\medskip\noindent
By Groupoids, Lemma \href{https://stacks.math.columbia.edu/tag/03BJ}{lemma integral over invariants (Tag 03BJ)}
the composition
$$
U \to \mathcal{X} \to M
$$
is an integral morphism of affine schemes. In particular, it is
universally closed
(Morphisms, Lemma \href{https://stacks.math.columbia.edu/tag/01WM}{lemma integral universally closed (Tag 01WM)}).
Since $U \to \mathcal{X}$ is surjective, it follows that $\mathcal{X} \to M$
is universally closed (Morphisms of Stacks, Lemma
\href{https://stacks.math.columbia.edu/tag/0CQK}{lemma image proper is proper (Tag 0CQK)}).
To conclude that $\mathcal{X} \to M$ is a universal homeomorphism,
it is enough to show that it is universally bijective, i.e.,
surjective and universally injective.

\medskip\noindent
We have $|\mathcal{X}| = |U|/|R|$ by
Morphisms of Stacks, Lemma \href{https://stacks.math.columbia.edu/tag/0DTZ}{lemma points presentation (Tag 0DTZ)}.
Thus $|f|$ is surjective and even bijective
by Groupoids, Lemma \href{https://stacks.math.columbia.edu/tag/03BL}{lemma points (Tag 03BL)}.

\medskip\noindent
Let $C \to C'$ be a ring map. Let $(U', R', s', t', c')$ be
the base change of $(U, R, s, t, c)$ by $M' = \Spec(C') \to M$.
Setting $\mathcal{X}' = [U'/R']$, we observe that
$M' \times_M \mathcal{X} = \mathcal{X}'$ by
Quotients of Groupoids, Lemma
\href{https://stacks.math.columbia.edu/tag/0DTG}{lemma base change quotient stack (Tag 0DTG)}.
Let $C^1$ be the ring of $R'$-invariant functions on $U'$.
Set $M^1 = \Spec(C^1)$ and consider the diagram
$$
\xymatrix{
\mathcal{X}' \ar[d]^{f'} \ar[r] & \mathcal{X} \ar[dd]^f \\
M^1 \ar[d] \\
M' \ar[r] & M
}
$$
By Groupoids, Lemma \href{https://stacks.math.columbia.edu/tag/03BK}{lemma invariants base change (Tag 03BK)} and
Algebra, Lemma \href{https://stacks.math.columbia.edu/tag/09EF}{lemma universally bijective (Tag 09EF)}
the morphism $M^1 \to M'$ is a homeomorphism.
On the other hand, the previous paragraph applied to
$(U', R', s', t', c')$ shows that $|f'|$ is bijective.
We conclude that $f$ induces a bijection on points after any
base change by an affine scheme. Thus $f$ is universally injective
by Morphisms of Stacks, Lemma
\href{https://stacks.math.columbia.edu/tag/0DTN}{lemma universally injective local (Tag 0DTN)}.

\medskip\noindent
Finally, we still have to show that $f$ is a uniform moduli space
in the category of affine schemes. This follows from the discussion
above and the fact that if the
ring map $C \to C'$ is flat, then $C' \to C^1$ is an isomorphism
by Groupoids, Lemma \href{https://stacks.math.columbia.edu/tag/03BK}{lemma invariants base change (Tag 03BK)}.
\end{proof}


\begin{lemma}
\label{lemma-etale-separated-over-well-nigh-affine}
Let $h : \mathcal{X}' \to \mathcal{X}$ be a morphism of algebraic stacks.
Assume $\mathcal{X}$ is well-nigh affine, $h$ is \'etale, $h$ is separated,
and $h$ induces isomorphisms on automorphism groups
(Morphisms of Stacks, Remark
\href{https://stacks.math.columbia.edu/tag/0DTW}{remark identify automorphism groups (Tag 0DTW)}).
Then there exists a cartesian diagram
$$
\xymatrix{
\mathcal{X}' \ar[d] \ar[r] & \mathcal{X} \ar[d] \\
M' \ar[r] & M
}
$$
where $M' \to M$ is a separated \'etale morphism of schemes and
$\mathcal{X} \to M$ is the moduli space constructed in
Lemma \ref{lemma-well-nigh-affine-moduli-space}.
\end{lemma}

\begin{proof}
Choose an affine scheme $U$ and a surjective, flat, finite, and
locally finitely presented morphism $U \to \mathcal{X}$.
Since $h$ is representable by algebraic spaces
(Morphisms of Stacks, Lemmas
\href{https://stacks.math.columbia.edu/tag/0DU9}{lemma aut iso unramified (Tag 0DU9)} and
\href{https://stacks.math.columbia.edu/tag/0DU7}{lemma stabilizer preserving (Tag 0DU7)})
we see that $U' = \mathcal{X}' \times_\mathcal{X} U$ is
an algebraic space. Since $U' \to U$ is separated and \'etale,
we see that $U'$ is a scheme and that every finite set of points
of $U'$ is contained in an affine open, see
Morphisms of Spaces, Lemma
\href{https://stacks.math.columbia.edu/tag/0418}{lemma locally quasi finite separated representable (Tag 0418)}
and
More on Morphisms, Lemma
\href{https://stacks.math.columbia.edu/tag/07S0}{lemma separated locally quasi finite over affine (Tag 07S0)}.
Setting $R' = U' \times_{\mathcal{X}'} U'$ we see
that $s', t' : R' \to U'$ are finite locally free.
By Groupoids, Lemma \href{https://stacks.math.columbia.edu/tag/03JE}{lemma find invariant affine (Tag 03JE)}
there exists an open covering $U' = \bigcup U'_i$ by
$R'$-invariant affine open subschemes $U'_i \subset U'$.
Let $\mathcal{X}'_i \subset \mathcal{X}'$ be the corresponding
open substacks. These are well-nigh affine as $U'_i \to \mathcal{X}'_i$
is surjective, flat, finite and of finite presentation. By
Lemma \href{https://stacks.math.columbia.edu/tag/0DUQ}{lemma well nigh affine moduli space etale (Tag 0DUQ)}
we obtain cartesian diagrams
$$
\xymatrix{
\mathcal{X}'_i \ar[r] \ar[d] & \mathcal{X} \ar[d] \\
M'_i \ar[r] & M
}
$$
with $M'_i \to M$ an \'etale morphism of affine schemes
and vertical arrows as in
Lemma \ref{lemma-well-nigh-affine-moduli-space}.
Observe that
$\mathcal{X}'_{ij} = \mathcal{X}'_i \cap \mathcal{X}'_j$
is an open subspace of $\mathcal{X}'_i$ and $\mathcal{X}'_j$.
Hence we get corresponding open subschemes
$V_{ij} \subset M'_i$ and $V_{ji} \subset M'_j$.
By the result of
Lemma \href{https://stacks.math.columbia.edu/tag/0DUR}{lemma moduli space finite affine (Tag 0DUR)}
we see that both
$\mathcal{X}'_{ij} \to V_{ij}$ and
$\mathcal{X}'_{ji} \to V_{ji}$ are categorical moduli spaces!
Thus we get a unique isomorphism $\varphi_{ij} : V_{ij} \to V_{ji}$
such that
$$
\xymatrix{
\mathcal{X}'_i \ar[d] & &
\mathcal{X}'_i \cap \mathcal{X}'_j \ar[rr] \ar[ll] \ar[ld] \ar[rd] & &
\mathcal{X}'_j \ar[d] \\
M'_i &
V_{ij} \ar[l] \ar[rr]^{\varphi_{ij}} & &
V_{ji} \ar[r] &
M'_j
}
$$
is commutative. These isomorphisms satisfy the cocyclce condition of
Schemes, Section \href{https://stacks.math.columbia.edu/tag/01JA}{section glueing schemes (Tag 01JA)} by a computation
(and another application of the previous lemma) which we omit.
Thus we can glue the affine schemes in to scheme $M'$, see
Schemes, Lemma \href{https://stacks.math.columbia.edu/tag/01JB}{lemma glue (Tag 01JB)}.
Let us identify the $M'_i$ with their image in $M'$.
We claim there is a morphism $\mathcal{X}' \to M'$ fitting into
cartesian diagrams
$$
\xymatrix{
\mathcal{X}'_i \ar[r] \ar[d] & \mathcal{X}' \ar[d] \\
M'_i \ar[r] & M'
}
$$
This is clear from the description of the morphisms into the glued scheme $M'$
in Schemes, Lemma \href{https://stacks.math.columbia.edu/tag/01JB}{lemma glue (Tag 01JB)} and the fact that to give a morphism
$\mathcal{X}' \to M'$ is the same thing as given a morphism $T \to M'$
for any morphism $T \to \mathcal{X}'$.
Similarly, there is a morphism $M' \to M$ restricting to the
given morphisms $M'_i \to M$ on $M'_i$.
The morphism $M' \to M$ is \'etale (being \'etale on the members of an
\'etale covering) and the fibre product property holds as it can
be checked on members of the (affine) open covering $M' = \bigcup M'_i$.
Finally, $M' \to M$ is separated because the composition
$U' \to \mathcal{X}' \to M'$ is surjective and universally closed
and we can apply Morphisms, Lemma
\href{https://stacks.math.columbia.edu/tag/09MQ}{lemma image universally closed separated (Tag 09MQ)}.
\end{proof}


\section*{Primary problem: Tag 0DUT}
\begin{theorem}[Keel-Mori]
\label{theorem-keel-mori}
Let $\mathcal{X}$ be an algebraic stack. Assume
$\mathcal{I}_\mathcal{X} \to \mathcal{X}$ is finite.
Then there exists a uniform categorical moduli space
$$
f : \mathcal{X} \longrightarrow M
$$
and $f$ is separated, quasi-compact, and a universal homeomorphism.
\end{theorem}

\begin{proof}
We choose a set $I$\footnote{The reader who is still keeping
track of set theoretic issues should make sure $I$ is not too large.}
and for $i \in I$ a morphism of algebraic stacks
$g_i : \mathcal{X}_i \to \mathcal{X}$ as in
Lemma \ref{lemma-etale-local-finite-inertia}; we will use all
of the properties listed in this lemma without further mention.
Let
$$
f_i : \mathcal{X}_i \to M_i
$$
be as in Lemma \ref{lemma-well-nigh-affine-moduli-space}.
Consider the stacks
$$
\mathcal{X}_{ij} = \mathcal{X}_i \times_{g_i, \mathcal{X}, g_j} \mathcal{X}_j
$$
for $i, j \in I$. The projections $\mathcal{X}_{ij} \to \mathcal{X}_i$
and $\mathcal{X}_{ij} \to \mathcal{X}_j$ are separated
by Morphisms of Stacks, Lemma
\href{https://stacks.math.columbia.edu/tag/050F}{lemma base change separated (Tag 050F)},
\'etale by Morphisms of Stacks, Lemma
\href{https://stacks.math.columbia.edu/tag/0CIN}{lemma base change etale (Tag 0CIN)},
and induce isomorphisms on automorphism groups
(as in Morphisms of Stacks, Remark
\href{https://stacks.math.columbia.edu/tag/0DTW}{remark identify automorphism groups (Tag 0DTW)}) by
Morphisms of Stacks, Lemma
\href{https://stacks.math.columbia.edu/tag/0DUB}{lemma base change stabilizer preserving (Tag 0DUB)}.
Thus we may apply Lemma \ref{lemma-etale-separated-over-well-nigh-affine}
to find a commutative diagram
$$
\xymatrix{
\mathcal{X}_i \ar[d]_{f_i} &
\mathcal{X}_{ij} \ar[d]_{f_{ij}} \ar[l] \ar[r] &
\mathcal{X}_j \ar[d]_{f_j} \\
M_i &
M_{ij} \ar[l] \ar[r] &
M_j
}
$$
with cartesian squares where $M_{ij} \to M_i$ and $M_{ij} \to M_j$
are separated \'etale morphisms of schemes; here we also use that $f_i$
is a uniform categorical quotient by
Lemma \href{https://stacks.math.columbia.edu/tag/0DUR}{lemma moduli space finite affine (Tag 0DUR)}.
Claim:
$$
\coprod M_{ij} \longrightarrow \coprod M_i \times \coprod M_i
$$
is an \'etale equivalence relation.

\medskip\noindent
Proof of the claim. Set $R = \coprod M_{ij}$ and $U = \coprod M_i$.
We have already seen that $t : R \to U$ and $s : R \to U$ are \'etale.
Let us construct a morphism $c : R \times_{s, U, t} R \to R$
compatible with $\text{pr}_{13} : U \times U \times U \to U \times U$.
Namely, for $i, j, k \in I$ we consider
$$
\mathcal{X}_{ijk} =
\mathcal{X}_i \times_{g_i, \mathcal{X}, g_j} \mathcal{X}_j
\times_{g_j, \mathcal{X}, g_k} \mathcal{X}_k =
\mathcal{X}_{ij} \times_{\mathcal{X}_j} \mathcal{X}_{jk}
$$
Arguing exactly as in the previous paragraph,
we find that $M_{ijk} = M_{ij} \times_{M_j} M_{jk}$
is a categorical moduli space for $\mathcal{X}_{ijk}$.
In particular, there is a canonical morphism
$M_{ijk} = M_{ij} \times_{M_j} M_{jk} \to M_{ik}$
coming from the projection $\mathcal{X}_{ijk} \to \mathcal{X}_{ik}$.
Putting these morphisms together we obtain the morphism $c$.
In a similar fashion we construct a morphism $e : U \to R$
compatible with $\Delta : U \to U \times U$ and
$i : R \to R$ compatible with the flip $U \times U \to U \times U$.
Let $k$ be an algebraically closed field. Then
$$
\Mor(\Spec(k), \mathcal{X}_i) \to \Mor(\Spec(k), M_i) = M_i(k)
$$
is bijective on isomorphism classes and the same remains true after any
base change by a morphism $M' \to M$. This follows from our choice
of $f_i$ and Morphisms of Stacks, Lemmas
\href{https://stacks.math.columbia.edu/tag/0CPN}{lemma universally injective (Tag 0CPN)} and
\href{https://stacks.math.columbia.edu/tag/0DTM}{lemma universally injective point (Tag 0DTM)}.
By construction of $2$-fibred products the diagram
$$
\xymatrix{
\Mor(\Spec(k), \mathcal{X}_{ij}) \ar[d] \ar[r] &
\Mor(\Spec(k), \mathcal{X}_j) \ar[d] \\
\Mor(\Spec(k), \mathcal{X}_i) \ar[r] &
\Mor(\Spec(k), \mathcal{X})
}
$$
is a fibre product of categories. By our choice of $g_i$ the
functors in this diagram induce bijections on automorphism groups.
It follows that this diagram induces a fibre product diagram
on sets of isomorphism classes! Thus we see that
$$
R(k) = U(k) \times_{|\Mor(\Spec(k), \mathcal{X})|} U(k)
$$
where $|\Mor(\Spec(k), \mathcal{X})|$ denotes the set
of isomorphism classes.
In particular, for any algebraically closed field $k$
the map on $k$-valued point is an equivalence relation.
We conclude the claim holds by
Groupoids, Lemma \href{https://stacks.math.columbia.edu/tag/0DT7}{lemma etale equivalence relation (Tag 0DT7)}.

\medskip\noindent
Let $M = U/R$ be the algebraic space which is the quotient of the above
\'etale equivalence relation, see
Spaces, Theorem \href{https://stacks.math.columbia.edu/tag/02WW}{theorem presentation (Tag 02WW)}.
There is a canonical morphism $f : \mathcal{X} \to M$
fitting into commutative diagrams
\begin{equation}
\label{equation-fundamental-diagram}
\xymatrix{
\mathcal{X}_i \ar[r]_{g_i} \ar[d]_{f_i} & \mathcal{X} \ar[d]^f \\
M_i \ar[r] & M
}
\end{equation}
Namely, such a morphism $f$ is given by a functor
$$
f : \Mor(T, \mathcal{X}) \longrightarrow \Mor(T, M)
$$
for any scheme $T$ compatible with base change. Let $a : T \to \mathcal{X}$
be an object of the left hand side. We obtain an \'etale covering
$\{T_i \to T\}$ with $T_i = \mathcal{X}_i \times_\mathcal{X} T$
and morphisms $a_i : T_i \to \mathcal{X}_i$. Then we get
$b_i = f_i \circ a_i : T_i \to M_i$. Since
$T_i \times_T T_j = \mathcal{X}_{ij} \times_\mathcal{X} T$
we moreover get a morphism $a_{ij} : T_i \times_T T_j \to \mathcal{X}_{ij}$.
Setting $b_{ij} = f_{ij} \circ a_{ij}$ we find that
$b_i \times b_j$ factors through the monomorphism
$M_{ij} \to M_i \times M_j$. Hence the morphisms
$$
T_i \xrightarrow{b_i} M_i \to M
$$
agree on $T_i \times_T T_j$. As $M$ is a sheaf for the \'etale
topology, we see that these morphisms glue to a unique
morphism $b = f(a) : T \to M$. We omit the verification that
this construction is compatible with base change and we omit
the verification that the diagrams
(\ref{equation-fundamental-diagram}) commute.

\medskip\noindent
Claim: the diagrams (\ref{equation-fundamental-diagram}) are cartesian.
To see this we study the induced morphism
$$
h_i : \mathcal{X}_i \longrightarrow M_i \times_M \mathcal{X}
$$
This is a morphism of stacks \'etale over $\mathcal{X}$
and hence $h_i$ is \'etale (Morphisms of Stacks, Lemma
\href{https://stacks.math.columbia.edu/tag/0CIR}{lemma etale permanence (Tag 0CIR)}).
Since $g_i$ is separated, we see $h_i$ is separated
(use Morphisms of Stacks, Lemma
\href{https://stacks.math.columbia.edu/tag/050M}{lemma compose after separated (Tag 050M)} and the fact
seen above that the diagonal of $\mathcal{X}$ is separated).
The morphism $h_i$ induces isomorphisms on automorphism groups
(Morphisms of Stacks, Remark
\href{https://stacks.math.columbia.edu/tag/0DTW}{remark identify automorphism groups (Tag 0DTW)})
as this is true for $g_i$. For an algebraically closed field $k$
the diagram
$$
\xymatrix{
\Mor(\Spec(k), M_i \times_M \mathcal{X}) \ar[r] \ar[d] &
\Mor(\Spec(k), \mathcal{X}) \ar[d] \\
M_i(k) \ar[r] &
M(k)
}
$$
is a catesian diagram of categories and the top arrow
induces bijections on automorphism groups.
On the other hand, we have
$$
M(k) = U(k)/R(k) = U(k)/
U(k) \times_{|\Mor(\Spec(k), \mathcal{X})|} U(k) =
|\Mor(\Spec(k), \mathcal{X})|
$$
by what we said above. Thus the right vertical arrow in the
cartesian diagram above is a bijection on isomorphism classes.
We conclude that
$|\Mor(\Spec(k), M_i \times_M \mathcal{X})| \to M_i(k)$ is bijective.
Review: $h_i$ is a separated, \'etale, induces isomorphisms on
automorphism groups (as in Morphisms of Stacks, Remark
\href{https://stacks.math.columbia.edu/tag/0DTW}{remark identify automorphism groups (Tag 0DTW)}), and
induces an equivalence on fibre categories over algebraically closed fields.
Hence it is an isomorphism by Morphisms of Stacks, Lemma
\href{https://stacks.math.columbia.edu/tag/0DUD}{lemma etale iso (Tag 0DUD)}.

\medskip\noindent
From the claim we get in particular the following:
we have a surjective \'etale morphism $U \to M$
such that the base change of $f$ is separated, quasi-compact,
and a universal homeomorphism. It follows that $f$ is separated,
quasi-compact, and a universal homeomorphism.
See Morphisms of Stacks, Lemma
\href{https://stacks.math.columbia.edu/tag/06TZ}{lemma check separated covering (Tag 06TZ)},
\href{https://stacks.math.columbia.edu/tag/0DTL}{lemma check quasi compact covering (Tag 0DTL)}, and
\href{https://stacks.math.columbia.edu/tag/0DTQ}{lemma check universal homeomorphism covering (Tag 0DTQ)}

\medskip\noindent
To finish the proof we have to show that $f : \mathcal{X} \to M$
is a uniform categorical moduli space.
To prove this it suffices to show that given a flat morphism
$M' \to M$ of algebraic spaces, the base change
$$
M' \times_M \mathcal{X} \longrightarrow M'
$$
is a categorical moduli space. Thus we consider a morphism
$$
\theta : M' \times_M \mathcal{X} \longrightarrow E
$$
where $E$ is an algebraic space. For each $i$ we know that
$f_i$ is a uniform categorical moduli space. Hence we obtain
$$
\xymatrix{
M' \times_M \mathcal{X}_i \ar[d] \ar[r] &
M' \times_M \mathcal{X} \ar[d]^\theta \\
M' \times_M M_i \ar[r]^{\psi_i} &
E
}
$$
Since $\{M' \times_M M_i \to M'\}$ is an \'etale covering,
to obtain the desired morphism $\psi : M' \to E$ it suffices
to show that $\psi_i$ and $\psi_j$ agree over
$M' \times_M M_i \times_M M_j = M' \times_M M_{ij}$.
This follows easily from the fact that
$f_{ij} : \mathcal{X}_{ij} =
\mathcal{X}_i \times_\mathcal{X} \mathcal{X}_j \to M_{ij}$ is a uniform
categorical quotient; details omitted.
Then finally one shows that $\psi$ fits into the commutative diagram
$$
\xymatrix{
M' \times_M \mathcal{X} \ar[d] \ar[rd]^\theta \\
M' \ar[r]^\psi &
E
}
$$
because ``$\{M' \times_M \mathcal{X}_i \to M' \times_M \mathcal{X}\}$
is an \'etale covering'' and the morphisms $\psi_i$ fit into the
corresponding commutative diagrams by construction.
This finishes the proof of the Keel-Mori theorem.
\end{proof}
\end{document}
